% s02-4-proof.tex —— 2.3 节（二）：命题 2.3.2.1 的四条（印刷页 45–54）

\subsection{\texorpdfstring{映射 $\Bun^{I_M,(S)}_M\to {}^{w}\Bun^{I,(S)}_P\times_{\Bun^I_G}{}^{w}\Bun^{I,(S)}_{P^-}$ 是开嵌入}{映射 IM,(S)M wI,(S)PIGwI,(S)P- 是开嵌入}}
\label{sec:four-a}

本小节证明命题~\ref{prop:four-conditions} 的 (a)，即
\[
\Bun^{I,(S)}_M\to
{}^{w}\Bun^{I,(S)}_P\times_{\Bun^I_G}{}^{w}\Bun^{I,(S)}_{P^-}
\]
是开嵌入。

这立刻由命题~\ref{prop:square-M-P} 的 (a)——即
$\Bun^{(S)}_M\to \Bun^{(S)}_P\times_{\Bun_G}\Bun^{(S)}_{P^-}$ 是开嵌入——结合下面
的命题得出。

\begin{proposition}
\label{prop:open-embedding-iso}
映射
\[
\Bun^{I,(S)}_M\to
\Bun^{(S)}_M\times_{\Bun^{(S)}_P\times_{\Bun_G}\Bun^{(S)}_{P^-}}
\bigl({}^{w}\Bun^{I,(S)}_P\times_{\Bun^I_G}{}^{w}\Bun^{I,(S)}_{P^-}\bigr)
\]
是同构。
\end{proposition}

\begin{proof}
注意展开左端得到
\[
\Bun^{(S)}_M\times_{\pt/M}\pt/\uB .
\]
对右端，先注意
\[
\begin{tikzcd}[column sep=large,row sep=large]
\Bun^{(S)}_M\times_{\pt/G}\pt/\uB \arrow[r] \arrow[d] \arrow[dr,phantom,"\lrcorner" description]
& \Bun^{I,(S)}_P\times_{\Bun^I_G}\Bun^{I,(S)}_{P^-} \arrow[d]\\
\Bun^{(S)}_M \arrow[r] & \Bun^{(S)}_P\times_{\Bun_G}\Bun^{(S)}_{P^-},
\end{tikzcd}
\]
\srcpage{45}
不过 $\Bun^{I,(S)}_P\times_{\Bun^I_G}\Bun^{I,(S)}_{P^-}$ 又可沿方
$(\Bun^{(S)}_P\times_{\Bun_G}\Bun^{(S)}_{P^-})\leftarrow \Bun^{(S)}_M$ 拉回，
其纤维是 $P\backslash G/B\times P^-\backslash G/B$；于是
\[
\begin{tikzcd}[column sep=large,row sep=large]
{}^{w}\Bun^{I,(S)}_P\times_{\Bun^I_G}{}^{w}\Bun^{I,(S)}_{P^-} \arrow[r] \arrow[d]
\arrow[dr,phantom,"\lrcorner" description]
& P\backslash PwB/B\times P^-\backslash P^-wB/B \arrow[d]\\
\Bun^{I,(S)}_P\times_{\Bun^I_G}\Bun^{I,(S)}_{P^-} \arrow[r]
& P\backslash G/B\times P^-\backslash G/B,
\end{tikzcd}
\]
\srcpage{46}
而 $\Bun^{I,(S)}_P\times_{\Bun^I_G}\Bun^{I,(S)}_{P^-}$ 沿
$\Bun^{(S)}_M\to \Bun^{(S)}_P\times_{\Bun_G}\Bun^{(S)}_{P^-}$ 拉回即得
$\Bun^{(S)}_M\times_{\pt/G}\pt/\uB$。因此展开右端得到
\[
\bigl(\Bun^{(S)}_M\times_{\pt/G}\pt/\uB\bigr)
\times_{P\backslash G/B\times P^-\backslash G/B}
\bigl(P\backslash PwB/B\times P^-\backslash P^-wB/B\bigr).
\]
而映射
\[
\Bun^{(S)}_M\times_{\pt/M}\pt/\uB\to
\bigl(\Bun^{(S)}_M\times_{\pt/G}\pt/\uB\bigr)
\times_{P\backslash G/B\times P^-\backslash G/B}
\bigl(P\backslash PwB/B\times P^-\backslash P^-wB/B\bigr)
\]
是沿 $\Bun^{(S)}_M\to \pt/M$ 对同构
\begin{align}
\pt/\uB&\to M\backslash G/B
\times_{P\backslash G/B\times P^-\backslash G/B}
\bigl(P\backslash PwB/B\times P^-\backslash P^-wB/B\bigr)
\label{eq:2.3.1}\\
&=M\backslash (PwB\cap P^-wB)/B
\label{eq:2.3.2}\\
&=M\backslash MwB/B
\label{eq:2.3.3}
\end{align}
\srcpage{47}
作基变换。
\end{proof}

\subsection{\texorpdfstring{映射 $\Bun^{I_M,(S)}_M\to {}^{w}\Bun^{I,(S)}_P$ 满}{映射 IM,(S)M wI,(S)P 满}}
\label{sec:four-b}

本小节证明命题~\ref{prop:four-conditions} 的 (b)，即对 Iwahori-收缩 $P$-可容子集 $S$，
\[
\Bun^{I_M,(S)}_M\to {}^{w}\Bun^{I,(S)}_P
\]
满。下面的图使论证更易跟随：
\[
\begin{tikzcd}[column sep=large,row sep=large]
\Bun^{I_M,(S)}_M \arrow[r] \arrow[d] &
{}^{w}\Bun^{I,(S)}_P\times_{\Bun^I_G}{}^{w}\Bun^{I,(S)}_{P^-} \arrow[r] \arrow[d] &
{}^{w}\Bun^{I,(S)}_{P^-} \arrow[d]\\
\Bun^{(S)}_M \arrow[r] &
\Bun^{(S)}_P\times_{\Bun_G}\Bun^{(S)}_{P^-} \arrow[r] &
\Bun^{(S)}_{P^-}.
\end{tikzcd}
\]

\begin{proof}
设 $y\in {}^{w}\Bun^{I,(S)}_P(k)$ 为 $k$-点，$F_P\in \Bun^{(S)}_P(k)$ 为其下方的
$P$-丛。由命题~\ref{prop:square-M-P} 的 (b)——即 $\Bun^{(S)}_M\to \Bun^{(S)}_P$ 满——
存在 $F_M\in \Bun^{(S)}_M(k)$ 使得 $F_P=F_M\times^M P$。考虑 $F_M$ 在映射
\[
\Bun^{(S)}_M\to \Bun^{(S)}_P\times_{\Bun_G}\Bun^{(S)}_{P^-}
\]
下的像。它与 $y$ 合起来给出纤维积
\begin{align*}
{}^{w}\Bun^{I,(S)}_P\times_{\Bun^{(S)}_P}
\bigl(\Bun^{(S)}_P\times_{\Bun_G}\Bun^{(S)}_{P^-}\bigr)
&={}^{w}\Bun^{I,(S)}_P\times_{\Bun^I_G}\Bun^{I,(S)}_{P^-}\\
&\supset {}^{w}\Bun^{I,(S)}_P\times_{\Bun^I_G}{}^{w}\Bun^{I,(S)}_{P^-}
\end{align*}
的一个 $k$-点。

具体地说，可把这个 $k$-点看作数据
\[
\bigl(F_P\in \Bun^{(S)}_P(k),\ F_{P^-}\in \Bun^{(S)}_{P^-}(k),\ F_{B,x}\in \pt/B(k)\bigr)
\]
连同同构
\[
(\gamma\colon F_P\times^P G\xrightarrow{\sim}F_{P^-}\times^{P^-}G;\quad
g\colon F_{B,x}\times^B G\xrightarrow{\sim}x^*F_P\times^P G),
\]
使得 $g\in (PwB)_{x^*F_M}\subset G_{x^*F_M}$，其中同构 $\gamma$ 与 $g$ 在 $x$ 处的
纤维分别给出 $G_{x^*F_M}$ 中的元素。取
$u\in U(P)_{x^*F_M}$ 使 $ug\in (MwB)_{x^*F_M}$。不妨把
$x^*\gamma\in G_{x^*F_M}$ 与单位元等同。于是注意到
\[
x^*\gamma\circ ug\in (MwB)_{x^*F_M}\subset (P^-wB)_{x^*F_M}.
\]
由推论~\ref{cor:UP-surj} 以及交换图
\[
\begin{tikzcd}[column sep=huge,row sep=large]
H^0(X,U(P)_{F_M}) \arrow[r] \arrow[d] & U(P)_{x^*F_M} \arrow[d]\\
H^0(X,G_{F_G}) \arrow[r] & G_{x^*F_G},
\end{tikzcd}
\]
\srcpage{49}
可知 $u$ 可提升为自同构 $\alpha\in \Aut(F_G)=H^0(G_{F_G})$。现考虑与前面主丛相同、
但同构改为
\[
(\gamma\circ \alpha,\ g)
\]
的 $k$-点。这样得到的 $k$-点落在
\[
{}^{w}\Bun^{I,(S)}_P\times_{\Bun^I_G}{}^{w}\Bun^{I,(S)}_{P^-}
\]
中。它与 $F_M$ 的数据合起来，由命题~\ref{prop:open-embedding-iso} 给出
$\Bun^{I,(S)}_M$ 的一个 $k$-点，该点在
\[
\Bun^{I,(S)}_M\to {}^{w}\Bun^{I,(S)}_P
\]
下的像正是 $y$。这正是所要证明的满性。
\end{proof}

\subsection{\texorpdfstring{映射 ${}^{w}\Bun^{I,(S)}_{P^-}\to \Bun^I_G$ 光滑}{映射 wI,(S)P- IG 光滑}}
\label{sec:four-c}

本小节证明命题~\ref{prop:four-conditions} 的 (c)，即对 Iwahori-收缩 $P$-可容子集 $S$，
\[
{}^{w}\Bun^{I,(S)}_{P^-}\to \Bun^I_G
\]
光滑。

由于该映射是光滑叠之间的局部有限表现态射，只需说明对点
\[
z\in {}^{w}\Bun^{I,(S)}_{P^-}
\]
有
\[
H^1\bigl(\mathrm{fib}(T_z\,{}^{w}\Bun^{I,(S)}_{P^-}\to T_{f(z)}\Bun^I_G)\bigr)=0 .
\]
\srcpage{50}
但 $S$ 是 $P$-可容的，故 ${}^{w}\Bun^{I,(S)}_{P^-}\subset {}^{w}\Bun^I_{P^-}$ 是开子叠，
于是只需说明
\[
H^1\bigl(\mathrm{fib}(T_z\,{}^{w}\Bun^I_{P^-}\to T_{f(z)}\Bun^I_G)\bigr)=0 .
\]
此处及以后记
\[
\mathrm{fib}=\mathrm{fib}\bigl(T_z\,{}^{w}\Bun^I_{P^-}\to T_{f(z)}\Bun^I_G\bigr).
\]
注意
\[
{}^{w}\Bun^I_{P^-}
=\Bun_{P^-}\times_{\pt/P^-}\pt/(B^w\cap P^-),
\]
故
\begin{align}
T_z\,{}^{w}\Bun^I_{P^-}
&=T_{F_{P^-}}\Bun_{P^-}\times_{T_{x^*F_{P^-}}\pt/P^-}
T_{F_{B^w\cap P^-,x}}\pt/(B^w\cap P^-)\label{eq:2.3.4}\\
&=\Gamma(X,\mathfrak{p}^-[1]_{F_{P^-}})\times_{\mathfrak{p}^-[1]x^*F_{P^-}}
(\mathfrak{b}^w\cap \mathfrak{p}^-)[1]_{F_{B^w\cap P^-,x}};\label{eq:2.3.5}
\end{align}
类似地，由于
\[
\Bun^I_G=\Bun_G\times_{\pt/G}\pt/B,
\]
有
\[
T_{f(z)}\Bun^I_G
=T_{F_G}\Bun_G\times_{T_{x^*F_G}\pt/G}T_{F_{B,x}}\pt/B
=\Gamma(X,\mathfrak{g}[1]_{F_G})\times_{\mathfrak{g}[1]x^*F_G}\mathfrak{b}[1]_{F_{B,x}},
\]
因此
\[
\mathrm{fib}
=\Gamma(X,(\mathfrak{g}/\mathfrak{p}^-)_{F_{P^-}})
\times_{(\mathfrak{g}/\mathfrak{p}^-)_{x^*F_{P^-}}}
\bigl(\mathfrak{b}/(\mathfrak{b}^w\cap \mathfrak{p}^-)\bigr)_{F_{B^w\cap P^-,x}},
\]
它给出长正合列
\[
H^0(X,(\mathfrak{g}/\mathfrak{p}^-)_{F_{P^-}})\oplus
\bigl(\mathfrak{b}/(\mathfrak{b}^w\cap\mathfrak{p}^-)\bigr)_{F_{B^w\cap P^-,x}}
\to (\mathfrak{g}/\mathfrak{p}^-)_{x^*F_{P^-}}\to H^1(\mathrm{fib})
\to H^1(X,(\mathfrak{g}/\mathfrak{p}^-)_{F_{P^-}}).
\]
\srcpage{51}
现在结论由如下观察得出：最后一项为 $0$，而第一个箭头由推论~\ref{cor:gp-global} 满。

\subsection{\texorpdfstring{映射 $\Bun^{I_M,(S)}_M\to {}^{w}\Bun^{I,(S)}_{P^-}$ 收缩}{映射 IM,(S)M wI,(S)P- 收缩}}
\label{sec:four-d}

本小节证明命题~\ref{prop:four-conditions} 的 (d)，即
\[
\Bun^{I_M,(S)}_M\to {}^{w}\Bun^{I,(S)}_{P^-}
\]
收缩。

策略是用命题~\ref{prop:contractive-criterion}。首先要写出一个映射
$\pi\colon {}^{w}\Bun^{I,(S)}_{P^-}\to \Bun^{I_M,(S)}_M$，使 $i$ 是它的截面。

注意有图
\[
\begin{tikzcd}[column sep=large,row sep=large]
{}^{w}\Bun^{I,(S)}_{P^-} \arrow[r] \arrow[d] & P^-\backslash P^-wB/B \arrow[d]\\
\Bun^{(S)}_{P^-} \arrow[r] & P^-\backslash G/B,
\end{tikzcd}
\]
其中两个方块均为拉回，且 $P^-\backslash P^-wB/B$ 可等同于 $\pt/(B^w\cap P^-)$，
于是有等同
\[
{}^{w}\Bun^{I,(S)}_{P^-}
=\Bun^{(S)}_{P^-}\times_{\pt/P^-}\pt/(B^w\cap P^-).
\]
在此等同下，映射 $i\colon \Bun^{I_M,(S)}_M\to {}^{w}\Bun^{I,(S)}_{P^-}$ 由图
\srcpage{52}
\[
\begin{tikzcd}[column sep=large,row sep=large]
\Bun^{I_M,(S)}_M \arrow[r] \arrow[d] & \pt/\uB \arrow[d]\\
\Bun^{(S)}_{P^-} \arrow[r] & \pt/(B^w\cap P^-),
\end{tikzcd}
\]
诱导；而映射 $\pi\colon {}^{w}\Bun^{I,(S)}_{P^-}\to \Bun^{I_M,(S)}_M$ 由图
\[
\begin{tikzcd}[column sep=large,row sep=large]
{}^{w}\Bun^{I,(S)}_{P^-} \arrow[r] \arrow[d] & \pt/(B^w\cap P^-) \arrow[d]\\
\Bun^{(S)}_M \arrow[r] & \pt/\uB
\end{tikzcd}
\]
诱导。于是很清楚 $i$ 是 $\pi$ 的截面。

下面在 $\Bun^{I_M,(S)}_M$ 上给 ${}^{w}\Bun^{I,(S)}_{P^-}$ 一个 $\mathbb{A}^1$-作用，
使得 $0$ 作用为 $i\circ\pi$，且诱导的 $\mathbb{G}_m$-作用在一个点上平凡。

仿 \cite[11.2 小节]{DG15}，取 $\mu\colon \mathbb{G}_m\to Z(M)\subset T$ 使得
$i\notin \Gamma_M$ 时 $\langle\mu,\check\alpha_i\rangle>0$。则 $P^-$ 上关于 $M$ 的
$\mathbb{G}_m$-作用
\[
\rho_t(p)=\mu(t)^{-1}p\mu(t)
\]
\srcpage{53}
延拓为 $P^-$ 上关于 $M$ 的 $\mathbb{A}^1$-作用，使得 $\rho_0\in \End(P^-)$ 对应于复合
$P^-\to M\to P^-$。此外，该作用把 $B^w\cap P^-\subset P^-$ 送到自身，于是得到
$B^w\cap P^-$ 上关于 $B^w\xrightarrow{\sim}B$ 的 $\mathbb{A}^1$-作用。
这给出 $\pt/P^-$ 上关于 $\pt/M$ 的 $\mathbb{A}^1$-作用，以及
$\pt/(B^w\cap P^-)$ 上与之相容的、关于 $\pt/\uB$ 的 $\mathbb{A}^1$-作用。
对 $\pt/P^-$ 上关于 $\pt/M$ 的 $\mathbb{A}^1$-作用施加 $\Maps(X,-)$，得到
$\Bun^{(S)}_{P^-}$ 上关于 $\Bun^{(S)}_M$ 的（自然相容的）$\mathbb{A}^1$-作用。

现映射
\[
{}^{w}\Bun^{I,(S)}_{P^-}\to \Bun^{I_M,(S)}_M
\]
形如
\[
\Bun^{(S)}_{P^-}\times_{\pt/P^-}\pt/(B^w\cap P^-)
\to \Bun^{(S)}_M\times_{\pt/M}\pt/\uB,
\]
而这些 $\mathbb{A}^1$-作用的纤维积满足命题~\ref{prop:contractive-criterion} 的假设，
正如所愿。
